\section{Perforce (Helix Core)}

\subsubsection*{p4d --- The Server}

\begin{tabularx}{\linewidth}{lX}
    \texttt{p4d}        & Start the server.\\
                        & \texttt{p4d -r <root> -p <port> -d}\\
    \texttt{-r <root>}  & Server root directory (\texttt{P4ROOT}):
                            holds metadata (\texttt{db.*}) and
                            versioned depot files.\\
    \texttt{-p <port>}  & Port to listen on, eg.\ \texttt{1666}.\\
    \texttt{-d}         & Run as a daemon (background).\\
\hline
\end{tabularx}

\begin{mdframed}[backgroundcolor=gray!10,linecolor=Firebrick4]
\begin{verbatim}
$ mkdir -p ~/perforce/root
$ p4d -r ~/perforce/root -p 1666 -d
\end{verbatim}
\end{mdframed}


\subsubsection*{Initial Admin Setup}

\begin{tabularx}{\linewidth}{lX}
    \texttt{p4 passwd}  & Set the admin (super user) password.
                            \textit{Run this first: with no
                            protections table yet, the connecting
                            user is implicitly super.}\\
\hline
    \texttt{p4admin}    & Admin console --- add users, create
                            depots, set protections, etc.\\
\hline
\end{tabularx}


\subsubsection*{Depots}

A \textbf{depot} is the top-level container for versioned files on
the server --- the closest equivalent to a Git ``repo''.

\begin{tabularx}{\linewidth}{lX}
    \texttt{local}  & \textit{Default} type. Plain versioned files,
                        addressed as \texttt{//depot/...}.\\
    \texttt{stream} & Files organized under Streams (managed
                        branching / merging).\\
\hline
\end{tabularx}
